\section{Avances}

% ── TABLA 1: Lujan De la Rosa · Cortez Vega · Arias Mendoza ──────────────────
\scriptsize
\begin{longtable}{|P{1.8cm}|P{4.4cm}|P{4.4cm}|P{4.4cm}|}
\hline
\rowcolor{colA}
{\color{white}\textbf{Fecha}} &
{\color{white}\textbf{Lujan De la Rosa Cristian Emir}} &
{\color{white}\textbf{Cortez Vega Jaime Fabian}} &
{\color{white}\textbf{Arias Mendoza Angel Adrian}} \\
\hline
\endhead

\rowcolor{colC}
12/10/2026 &
  Análisis de requerimientos, conformación de sub equipos y diseño de la arquitectura
  general. Definición de la estructura de las clases de modelo. &
  Análisis de requerimientos y estudio del algoritmo de clausura LR(0). Diseño de la
  función iterativa de punto fijo. &
  Análisis de requerimientos. Diseño de las clases \texttt{Production} y
  \texttt{Grammar} con sus validaciones. \\[3pt]
\hline

\rowcolor{colD}
13/10/2026 &
  Implementación del lector de gramáticas en formato posicional y pruebas con la
  Gramática~1. &
  Implementación de \texttt{closure} con iteración de punto fijo; primeras pruebas
  con la Gramática~1. &
  Implementación de \texttt{Production} y \texttt{Grammar} con \texttt{augmented()};
  pruebas de validaciones. \\[3pt]
\hline

\rowcolor{colC}
14/10/2026 &
  Extensión del lector al formato etiquetado (\texttt{inicio:}, \texttt{terminales:}).
  Soporte para flechas \texttt{->} y comentarios. &
  Implementación de \texttt{goto} e integración con \texttt{closure}. Pruebas con
  las Gramáticas~2 y~3. &
  Implementación de \texttt{LR0Item} con \texttt{advance()} y de \texttt{LR0State}
  con sus transiciones. \\[3pt]
\hline

\rowcolor{colD}
15/10/2026 &
  Pruebas del lector con las seis gramáticas de ejemplo e integración con el módulo
  de modelos. &
  Implementación de \texttt{construir\_coleccion} con generación de traza. Pruebas
  con las Gramáticas~4, 5 y~6. &
  Implementación de \texttt{CanonicalCollection}. Pruebas de integración con los
  algoritmos de clausura e ir\_a. \\[3pt]
\hline

\rowcolor{colC}
16/10/2026 &
  Pruebas finales del lector con archivos malformados; verificación de manejo de
  errores. Redacción del informe. &
  Pruebas finales y comparación de resultados con los esperados del profesor.
  Apoyo en el informe. &
  Pruebas finales del módulo de modelos. Redacción de conclusiones y revisión de
  la rúbrica. \\[3pt]
\hline

\end{longtable}

\vspace{0.4cm}
% ── TABLA 2: Vasquez Moreno · Nicolas Vasquez · Carrasco Pacheco ───────────────
\scriptsize
\begin{longtable}{|P{1.8cm}|P{4.4cm}|P{4.4cm}|P{4.4cm}|}
\hline
\rowcolor{colA}
{\color{white}\textbf{Fecha}} &
{\color{white}\textbf{Vasquez Moreno Zair De Jesus}} &
{\color{white}\textbf{Nicolas Vasquez Joel Osmar}} &
{\color{white}\textbf{Carrasco Pacheco Levi Josue}} \\
\hline
\endhead

\rowcolor{colC}
12/10/2026 &
  Análisis de requerimientos. Investigación manual del algoritmo LR(0) sobre
  gramáticas sencillas. &
  Análisis de requerimientos. Diseño del módulo de estilos y la paleta de colores. &
  Análisis de requerimientos. Diseño del componente \texttt{TablaEstados} con
  scroll vertical. \\[3pt]
\hline

\rowcolor{colD}
13/10/2026 &
  Práctica manual de clausura e ir\_a con la Gramática~1 para verificar resultados
  esperados. &
  Implementación de \texttt{estilos.py}: diccionario \texttt{COLORES} y temas
  \texttt{ttk} globales. &
  Implementación de \texttt{TablaEstados} con \texttt{Canvas} y barra de scroll
  vertical. \\[3pt]
\hline

\rowcolor{colC}
14/10/2026 &
  Apoyo en pruebas de \texttt{goto} con las Gramáticas~2 y~3. Verificación de
  estados generados. &
  Implementación de \texttt{VentanaResultados} con tabla de traza de dos columnas
  y doble scroll. &
  Implementación de \texttt{InterfazColeccion} con botones de carga, análisis y
  limpieza. \\[3pt]
\hline

\rowcolor{colD}
15/10/2026 &
  Pruebas del flujo completo con todas las gramáticas de ejemplo. Verificación del
  estado de aceptación. &
  Integración de \texttt{VentanaResultados} con la traza de
  \texttt{construir\_coleccion}; corrección de formato. &
  Pruebas de integración de la interfaz completa en macOS. Corrección de detalles
  visuales en la barra de estado. \\[3pt]
\hline

\rowcolor{colC}
16/10/2026 &
  Pruebas finales y verificación del número de estados con los resultados del
  profesor. Apoyo en el informe. &
  Redacción de los elementos del informe. Integración del diagrama UML en el
  documento. &
  Redacción final del documento, revisión de ortografía y verificación del PDF
  generado. \\[3pt]
\hline

\end{longtable}
\normalsize
